% Preambolo per i PDF legali/privacy — incluso da tools/legal/build_pdf.sh.
%
% Nasce da un problema concreto: nelle tabelle lunghe (registro art. 30
% Sezione D, matrice di rischio della DPIA) le righe si confondevano fra loro.
% Lo stile booktabs, usato da pandoc, per scelta non traccia linee fra le
% righe: va bene per tabelle compatte, meno quando una cella occupa piu' linee
% e il lettore perde la corrispondenza orizzontale.

% I comandi qui sotto vivono in booktabs (\heavyrulewidth), longtable (\LTpre)
% ed etoolbox (\AtBeginEnvironment). Il template di pandoc li carica soltanto
% dentro un $if(tables)$: un documento SENZA tabelle -- i ToS, per esempio --
% non li ha, e xelatex si ferma su "Undefined control sequence". Il PDF non
% viene prodotto e sul disco resta quello vecchio.
%
% Caricarli qui li rende disponibili in ogni caso; ripetere \usepackage con le
% stesse opzioni e' innocuo quando pandoc li ha gia' inclusi.
\usepackage{booktabs}
\usepackage{longtable}
\usepackage{etoolbox}

% Piu' aria sopra e sotto il testo di ogni cella: separa le righe senza
% aggiungere tratti, restando coerente con lo stile booktabs.
\renewcommand{\arraystretch}{1.35}

% Filetti leggermente piu' marcati in testa e in coda alla tabella, cosi' il
% blocco si stacca dal testo circostante.
\setlength{\heavyrulewidth}{1.1pt}
\setlength{\lightrulewidth}{0.5pt}

% Le tabelle larghe respirano meglio con un corpo leggermente ridotto: evita
% che una cella lunga si spezzi in troppe righe.
\AtBeginEnvironment{longtable}{\small}

% Spazio verticale attorno alla tabella.
\setlength{\LTpre}{10pt}
\setlength{\LTpost}{10pt}
